Si $(\xi_j)$ son diferencias de martingala con $|\xi_j|\le d_j$, para $\theta>0$ la convexidad de la
exponencial entre $-d_j$ y $d_j$ y la media condicional cero dan
$\E[\exp(\theta\xi_j)\mid\mathcal F_{j-1}]\le\cosh(\theta d_j)\le\exp(\theta^2d_j^2/2)$. El producto
compensado es una supermartingala positiva. Pararla en el primer $h$ con $\sum_{j\le h}\xi_j\ge u$,
aplicar la misma cuenta a $-\xi_j$ y optimizar $\theta$ da
\begin{equation}\label{macro:max-mart}
 P\left(\max_{h\le m}\left|\sum_{j\le h}\xi_j\right|\ge u\right)
 \le 2\exp\!\left(-\frac{u^2}{2\sum_{j\le m}d_j^2}\right).
\end{equation}
Es la forma maximal de la desigualdad de Azuma~\cite{azuma}. Para la partición $\mathcal I$ fija
existe $C'_{\mathcal I}$ tal que si todas las martingalas escalares
tienen máximo absoluto menor que $b/C'_{\mathcal I}$ entonces $\mathcal M^*<b$. El factor $C'_{\mathcal I}$
cuenta los intervalos de una sola fibra y las observaciones de intervalos y rectángulos. El número de
líneas distinguidas sólo interviene en la unión de sucesos. De \eqref{macro:max-mart} resulta
\[
 P^{\rm B}(\mathcal M^*\ge b)\le C_{\mathcal I}(n+1)\exp(-c_{\delta,\mathcal I}b^2 n).
\]

Al inicio $x=U_f=1$. Las sumas de retícula por intervalos y cajas, la aproximación de los extremos
de las fibras y las $O(n^3)$ cuádruplas excepcionales dan $\Delta(0)\le C_{\mathcal I}/n$,
uniformemente en la línea distinguida. Elegimos primero $\varrho_{\rm mac}$ dentro de los umbrales
deterministas, después $a_{\rm mac}$ y $b$ con
$e^{C_\delta}(b+C_\delta a_{\rm mac})<\varrho_{\rm mac}/3$, y finalmente $n$ grande a partición
fija. Gronwall discreto en \eqref{macro:gronwall} da $\Delta(s\wedge\tau)<\varrho_{\rm mac}$ siempre
que $\mathcal M^*<b$, también para el paso que llega a $\tau$. Así
$\mathcal H_M\subseteq\{\mathcal M^*\ge b\}$, lo que prueba la cota, absorbiendo $n+1$ en la
exponencial.

\emph{Refinamiento de la partición.} Queda la afirmación sobre una partición de observación fija.
Tomemos como $\mathcal I$ un refinamiento suyo de paso suficientemente pequeño. Usamos el árbol común
de órdenes de filas y prefijos de elecciones de
casillas disponibles, detenido en la primera salida de la partición fina o en el tamaño $M$. En cada
prefijo, las masas y los objetivos de la partición gruesa se obtienen sumando sus partes finas. El instante
de congelación de una línea depende sólo de las casillas elegidas, de modo que esto vale también para
las fibras congeladas. La desigualdad triangular, seguida para las fibras del máximo sobre las
líneas, da $\Delta_{\rm gruesa}\le\Delta_{\rm fina}$ en ese mismo prefijo, donde los subíndices
indican la partición con que se calcula la discrepancia. Por tanto, al radio común,
la partición gruesa no puede salir antes que la fina. Si un prefijo de salida de la partición fina
admite una continuación, toda salida posterior de la gruesa pertenece al suceso ya contado con la
masa total de ese prefijo. En este argumento sólo intervienen transiciones desde padres estrictamente
anteriores a la parada fina. Esos padres tienen punto de Bethe por la Proposición~\ref{macro:maximo}.
\end{proof}
